\scanpage{47}
À Algèbre de quaternions \inserted{$\mathcal M$}, on associe
\[
\left\{\begin{aligned}
 &\left.\begin{aligned}
 &\text{forme }G\text{ de }\operatorname{GP}(1)_S\\
 &\text{----- }\widetilde G\text{ --- }\operatorname{SL}(2)_S
 \end{aligned}\right\}\\[-2pt]
 &\left.\begin{aligned}
 &\text{----- }\mathfrak g\ \deleted{\text{de}}\text{ --- }\mathfrak{gp}(1)_S\\
 &\text{----- }\widetilde{\mathfrak g}\text{ --- }\unsure{\mathfrak{sl}|\mathfrak{gl}}(1)_S
 \end{aligned}\right\}\ \Big\downarrow\\
 &(E,q,\omega)\quad(E=\widetilde{\mathfrak g})\quad
 \text{\unsure{Module|module} spécial quadratique de rg 2}\\
 &P\supset X\qquad\text{quadrique dans plan projective (}\\
 &\underset{\mathbb P^\vee(E)}{\shortparallel}\quad
 \underset{\operatorname{Var}(q)}{\shortparallel}\\
 &X\text{ fibré en droites projectives}
\end{aligned}\right.
\]
On trouve que
\[
 \left\{\begin{aligned}
 &\mathcal M\mapsto G,\ \mathcal M\mathrel{\unsure{\mapsto|\vdash}}\widetilde G,\
 \mathcal M\mapsto(E,q,\omega),\ \text{et}\\
 &\mathcal M\deleted{\unsure{\ell}}\mapsto(P\supset X),\quad\mathcal M\mapsto X
 \end{aligned}\right.
\]
sont des équivalences de catégories \inserted{fibrées}, et que $\mathcal M\mapsto\mathfrak g,\mathcal M\mapsto\widetilde{\mathfrak g}$ sont des équiv. de cat. fibrées au dessus de $\operatorname{Sp}\mathbb Z(\tfrac12)$.

Considérons maintenant le cas où on a un fibré vect. de rang 2 $V$, et \deleted{\unsure{où on}}
\[
 \mathcal M=V^\vee\otimes V\simeq\operatorname{End}(V)\tag*{(1)}
\]
On va d'abord définir un iso can
\[
 X\xleftarrow[\sim]{\varphi}\mathbb P(V)
 \bigl(\overset{\inserted{\text{orth.}}}{\simeq}\mathbb P(V^\vee)\bigr)\tag*{(2)}
\]
Pour ceci, si \tikz[baseline=(l.base)]{\node[inner sep=1pt] (l) {$L\inserted{\subset V}$};\draw ([xshift=-1pt]l.south west) |- ([xshift=1pt,yshift=1pt]l.north east) -- ([xshift=1pt]l.east);\draw ([xshift=-3pt]l.south west) |- ([xshift=3pt,yshift=3pt]l.north east) -- ([xshift=3pt]l.east);} est un ss-fibré vectoriel inv. de $V$, \deleted{\unsure{$L'\ldots$}}
\[
 L'\subset V'\qquad L'=L^\perp
\]
son orthogonal, on leur associe
\[
 \varphi(L)=L'\otimes L=L^\perp\otimes L\subset V^\vee\otimes V=\mathcal M\tag*{(3)}
\]
et on trouve bien un iso \ldots
% archive sheet number: 46
